\chapter{联络、平行移动与曲率}
\label{ch:10-connections}
\ChapterPlan{周第3章（本地页码待核对）；梁上册第二版3.1--3.4及第3章习题1、2、6、8；一般向量丛与结构方程另作补充}{正文及证明初稿；附8题解答与原创图，尚未独立数学审阅}

不同基点的向量属于不同的纤维，不能直接相减。联络给出沿方向求导的规则，并通过微分方程比较不同纤维。本章从一般向量丛开始，以切丛联络为主线；不预先选度规，也不默认无挠。平行移动、测地线和曲率各有不同作用：平行移动沿给定曲线比较纤维，测地线要求速度自身平行，曲率则度量局部平行比较的不相容。

固定$\nabla_{\partial_i}\partial_j=\Gamma^k_{ij}\partial_k$，第一个下指标对应求导方向。曲率始终取$R(X,Y)=\nabla_X\nabla_Y-\nabla_Y\nabla_X-\nabla_{[X,Y]}$，与第2章及附录相同。梁书默认无挠、以余矢量的交换导数定义曲率分量，不能直接搬用其分量号差；本章会给出对偶作用的负号。

\Needspace{17\baselineskip}
\section{仿射联络与局部系数}
\begin{definition}[向量丛联络与仿射联络]\label{def:connection}
设$E\to M$为秩$r$的实光滑向量丛，记$\Gamma(E)$为光滑截面空间、$\mathfrak X(M)=\Gamma(TM)$。联络（connection）是实双线性映射
\[
 \nabla:\mathfrak X(M)\times\Gamma(E)\to\Gamma(E),\qquad (X,s)\mapsto\nabla_Xs,
\]
满足$\nabla_{fX}s=f\nabla_Xs$以及$\nabla_X(fs)=X(f)s+f\nabla_Xs$。$E=TM$时称仿射联络（affine connection）。$\nabla_Xs$称$s$沿$X$的协变导数（covariant derivative）。
\end{definition}
第一变量只依赖方向向量的点值，第二变量还需要截面的一阶变化。若第二变量也为$C^\infty$线性，Leibniz法则中的$X(f)s$就无从出现。抽象联络不必保持任何长度或内积，不能把“平行”先解释为等长。

\begin{proposition}[局部性与局部表达]\label{prop:connection-local}
联络可以限制到开子集，值$\nabla_Xs(p)$只依赖$X,s$在$p$附近的取值。若$e_1,\ldots,e_r$为局部标架，则有唯一矩阵值一形式$\omega=(\omega^a{}_b)$，使
\[
 \nabla_Xe_b=e_a\omega^a{}_b(X),\qquad
 \nabla_X(e_b s^b)=e_a\bigl(X(s^a)+\omega^a{}_b(X)s^b\bigr).
\]
反过来，任意光滑矩阵值一形式按此公式给出局部联络。
\end{proposition}
\begin{proof}
若$s$在$p$附近为零，选第4章的截断$\chi=1$于$p$附近且支集位于该邻域。于是$s=(1-\chi)s$，在$p$处用Leibniz法则得到$\nabla_Xs(p)=0$，因为$1-\chi$及其微分都在那里为零。对$X$在$p$附近为零的情形，同样用第一变量的函数线性性。将局部数据乘以这样的截断延成全局数据，所得点值与延拓无关，故可限制联络。

局部写$X=X^i\partial_i$，第一变量线性给$\nabla_Xe_b=X^i\nabla_{\partial_i}e_b$，所以各标架系数是$X$的点态线性函数，组成光滑一形式$\omega^a{}_b$。对$s=e_b s^b$应用Leibniz即得公式。反向直接逐项检验两条函数规则，且结果是光滑截面。
\end{proof}

\Needspace{17\baselineskip}
\begin{proposition}[换标架与坐标系数]\label{prop:connection-change}
将标架视为行向量$e=(e_1,\ldots,e_r)$，令$e'=eG$，其中$G:U\to\operatorname{GL}(r,\R)$光滑；截面系数列向量相应满足$s'=G^{-1}s$。则
\[
 \omega'=G^{-1}\omega G+G^{-1}\dd G.
\]
对切丛取坐标标架，$\omega^k{}_j=\Gamma^k_{ij}\dd x^i$，有
\begin{align*}
 (\nabla_XY)^k&=X^i\bigl(\partial_iY^k+\Gamma^k_{ij}Y^j\bigr),\\
 \Gamma'^a_{bc}&=\frac{\partial y^a}{\partial x^k}
 \left(\frac{\partial x^i}{\partial y^b}\frac{\partial x^j}{\partial y^c}\Gamma^k_{ij}
       +\frac{\partial^2x^k}{\partial y^b\partial y^c}\right).
\end{align*}
后一式是$y$坐标系数的变换律，故$\Gamma^k_{ij}$一般不是同一个张量在不同坐标下的分量。
\end{proposition}
\begin{proof}
对$e'_b=e_aG^a{}_b$求导得到$\nabla e'=e(\omega G+\dd G)=e'G^{-1}(\omega G+\dd G)$，从而得到矩阵公式。这里$\dd G$逐项求微分，矩阵乘法次序不能交换。

切丛局部式由上一命题直接得到。再写$\partial_{y^c}=(\partial x^j/\partial y^c)\partial_{x^j}$，沿$\partial_{y^b}$求导。导数打在系数上给二阶偏导，打在标架上给两个一次Jacobi因子与$\Gamma$；最后将输出$\partial_{x^k}$换成$(\partial y^a/\partial x^k)\partial_{y^a}$，即得所列式。
\end{proof}

\begin{theorem}[联络的存在及其差]\label{thm:connection-affine}
每个有限秩光滑向量丛都有联络。若$\nabla,\widehat\nabla$为同一丛上的两个联络，则
\[
 A(X)s=\widehat\nabla_Xs-\nabla_Xs
\]
由唯一的$A\in\Omega^1(M;\operatorname{End}E)$给出。反过来，任意这样的$A$都使$\nabla+A$成为联络。因此联络集合是以该向量空间为模型的仿射空间，未指定天然的零联络。
\end{theorem}
\begin{proof}
取局部平凡化覆盖，每块用$\omega_\alpha=0$定义联络$\nabla^\alpha$，以从属单位分解$\rho_i$令
\[
 \nabla_Xs=\sum_i\rho_i\nabla^{\alpha(i)}_X(s|_{U_{\alpha(i)}}),
\]
每项在域外延零。支集从属保证延零光滑，局部有限保证和光滑。第一变量线性立即成立；第二变量Leibniz的额外项是$\sum_i\rho_i X(f)s=X(f)s$，所以得到全局联络。

对差$A(X)(fs)$，两个Leibniz项$X(f)s$抵消，故它对$s$为函数线性；对$X$也函数线性。在局部标架下，这使其系数成为点态的双线性映射，给出光滑$T^*M\otimes\operatorname{End}E$截面。反向加上这样的零阶项不会改变联络规则。任意两个联络的差唯一，而仿射平均只有权重之和为$1$时才保留Leibniz系数。
\end{proof}

\begin{example}[极坐标中的非零系数]
平面标准联络在$(x,y)$坐标中$\Gamma=0$。取$r>0$的极坐标卡，直接将$\partial_r=(\cos\theta,\sin\theta)$、$\partial_\theta=(-r\sin\theta,r\cos\theta)$沿坐标方向求导，得到
\[
 \Gamma^r_{\theta\theta}=-r,\qquad
 \Gamma^\theta_{r\theta}=\Gamma^\theta_{\theta r}=1/r,
\]
其余为零。系数非零并不说明平面有曲率；在后面计算中，系数导数与二次项会抵消。把每张卡的$\Gamma$都设为零也不能直接定义同一个联络，因为换卡有二阶项。
\end{example}
\begin{figure}[htbp]
\centering\DGDiagram{connection-frame-change}

\caption{同一截面在两套标架中的系数和协变导数均按$G^{-1}$变换，局部算子为$D=\dd+\omega$。交换关系导出非齐次项$G^{-1}\dd G$。这里比较的是同一基点上的标架，非沿路径的平行移动。原创示意。}
\label{fig:connection-frame-change}
\end{figure}


\Needspace{17\baselineskip}
\section{平行移动与测地线方程}
\begin{definition}[沿曲线的协变导数]\label{def:derivative-along-curve}
设$\gamma:I\to M$光滑，$V(t)\in E_{\gamma(t)}$为沿曲线的光滑截面。在局部标架下写$V=e_a(\gamma(t))v^a(t)$，定义
\[
 D_tV=e_a(\gamma(t))\left(\dot v^a+
 \omega^a{}_b(\dot\gamma)v^b\right).
\]
若$D_tV=0$，称$V$沿$\gamma$平行（parallel）。对切丛亦写$D_tV=\nabla_{\dot\gamma}V$，但此式按沿曲线截面理解，不要求$V$来自某个邻域向量场。
\end{definition}
换标架时$v'=G(\gamma)^{-1}v$，利用上一节的$\omega'$公式和$\frac{\dd}{\dd t}G(\gamma)=\dd G(\dot\gamma)$，直接得到$\dot v'+\omega'(\dot\gamma)v'=G^{-1}(\dot v+\omega(\dot\gamma)v)$。因此定义无关标架。它适用于自交、驻点及常值曲线；在常值曲线上只剩固定纤维中向量的通常导数。若$V=s\circ\gamma$来自局部截面，则$D_tV=(\nabla_{\dot\gamma}s)\circ\gamma$，与原联络相容。

\begin{theorem}[平行移动的存在与基本性质]\label{thm:parallel-transport}
对紧区间$[a,b]$上的分段光滑曲线$\gamma$及$v\in E_{\gamma(a)}$，存在唯一连续、分段光滑的平行截面$V$，满足$V(a)=v$。由此定义线性同构
\[
 P^\gamma_{b\leftarrow a}:E_{\gamma(a)}\to E_{\gamma(b)},\qquad v\mapsto V(b).
\]
它满足恒等、拼接和逆向性质：
\[
 P^\gamma_{t\leftarrow s}P^\gamma_{s\leftarrow r}=P^\gamma_{t\leftarrow r},
 \qquad P^{\bar\gamma}_{b\leftarrow a}=(P^\gamma_{b\leftarrow a})^{-1},
\]
其中$\bar\gamma(t)=\gamma(a+b-t)$在同一个$[a,b]$上反向走；保持方向的光滑重参数化不改变移动。光滑有限维曲线族的平行移动局部光滑依赖于参数及初值。
\end{theorem}
\begin{proof}
单块中方程是$\dot v=-A(t)v$，$A(t)=\omega(\dot\gamma(t))$。第5章定理\ref{thm:ode-local}给局部存在唯一性与光滑参数依赖。若$\|A(t)\|\leq C$，Gronwall引理给$|v(t)|\leq|v(s)|e^{C|t-s|}$。在有限端点前解因而有界，且$|\dot v|\leq C|v|$有界，所以存在端点极限；以该极限为新初值可延拓，排除有限时刻爆破。故在线性方程系数定义的整个紧子区间都有解。

对一般曲线，紧时间区间可有限分割，使每小段的像位于一个平凡化域中；逐段解方程并在端点换标架。局部定义相容和唯一性保证选择分割不影响结果。解对初值线性，因为方程齐次且唯一；逆向解同一方程恢复原初值，所以移动可逆。沿相邻时间段拼接解，由唯一性得复合式。重参数化$\gamma(\phi(u))$时导数整体乘$\phi'(u)$，故平行方程保持，得到重参数化结论。

参数光滑依赖先在一个给定参数值附近取上述有限分割，缩小参数邻域后仍能使用同一组块，再逐段应用ODE参数定理。此陈述针对具有固定有限分段的有限维光滑曲线族，不在此建立无限维路径空间上的光滑结构。
\end{proof}
这里$P^\gamma_{t\leftarrow s}$的上下标为该路径参数的终点与起点；对反向路径，其源纤维为$E_{\gamma(b)}$、目标为$E_{\gamma(a)}$，所以确实是原移动的逆。
固定$p$，所有以$p$为起终点的分段光滑闭路移动组成$\operatorname{GL}(E_p)$的子群，称holonomy群（闭路平行移动群）。恒定路径给单位，路径拼接给乘积，反向给逆。若用一条从$p$到$q$的路径比较基点，两群通过该路径的移动共轭。因此“闭路移动非平凡”是联络的几何性质，不是所选标架造成的假象。
\begin{figure}[htbp]
\centering\DGDiagram{parallel-transport}

\caption{沿给定曲线的线性同构由平行方程决定。两条竖线示意不同基点的纤维，向量长度仅用于排版，不表示移动保持或改变某个已选度量；本章尚未选度规。原创示意。}
\label{fig:parallel-transport}
\end{figure}

\begin{example}[平行不等于分量恒定]
平面中恒定向量$V=\partial_x$沿半径$r=R$的圆周平行。其极坐标分量是$v^r=\cos\theta$、$v^\theta=-\sin\theta/R$，均随$\theta$变化。代入$D_\theta v^r=\partial_\theta v^r-Rv^\theta$及$D_\theta v^\theta=\partial_\theta v^\theta+v^r/R$，两者都为零。只有在沿该曲线的平行标架中，平行截面的系数才恒定。
\end{example}

\begin{definition}[仿射测地线]
切丛联络$\nabla$的仿射参数测地线（affinely parametrized geodesic）是满足$D_t\dot\gamma=0$的光滑曲线；局部方程为
\[
 \ddot x^k+\Gamma^k_{ij}(x)\dot x^i\dot x^j=0.
\]
此定义包括常值曲线；在尚无度规时，不把测地线解释为最短曲线。
\end{definition}

\begin{theorem}[测地线初值、喷流与参数]\label{thm:affine-geodesics}
给定$(p,v)\in TM$，存在唯一局部测地线满足$\gamma(0)=p,\dot\gamma(0)=v$，并光滑依赖初值；解可延拓为唯一极大开区间上的解。若$v\ne0$，其速度在定义域中处处非零。仿射重参数化$t=au+b$（$a\ne0$）保持方程；对非恒定测地线，正则重参数化保持仿射测地线方程当且仅当$\phi''=0$。测地线只依赖$\Gamma^k_{(ij)}$，不能单独确定挠率。
\end{theorem}
\begin{proof}
将方程写成$TM$中的一阶系统$\dot x^k=v^k$、$\dot v^k=-\Gamma^k_{ij}(x)v^iv^j$。对坐标变化$y=y(x)$，链式法则及系数变换律给
\[
 \ddot y^a+\Gamma'^a_{bc}\dot y^b\dot y^c
 =\frac{\partial y^a}{\partial x^k}
   \bigl(\ddot x^k+\Gamma^k_{ij}\dot x^i\dot x^j\bigr),
\]
其中$\ddot y$的坐标二阶项与$\Gamma'$中的逆坐标二阶项相消；这也可由恒等$y(x(y))=y$求两次导数验证。所以该一阶系统定义$TM$上的光滑向量场。第5章ODE与极大积分曲线理论给初值、光滑依赖和极大延拓。若某时刻速度为零，其后与常值曲线有相同初值，由唯一性反向追溯会使原解恒定，矛盾。

对$\widetilde\gamma(u)=\gamma(\phi(u))$，链式法则给
\[
 D_u\dot{\widetilde\gamma}=(\phi')^2D_t\dot\gamma+\phi''\dot\gamma.
\]
非恒定解的速度非零，所以右端在原解为测地线时为零恰要求$\phi''=0$。最后，反对称系数与$v^iv^j$缩并为零，方程只看到对称部分。
\end{proof}
\begin{example}[测地线不一定在所有时间存在]
在$M=(0,1)$取$\Gamma=0$，$x(t)=1/2+t$的极大区间是$(-1/2,1/2)$；联络系数光滑且曲率为零，并未阻止曲线离开流形。平行移动沿每条已经给定的紧时间路径都能完成，不表示测地线对任意时间均有定义。完备性、指数映射及长度变分留到第12章。
\end{example}


\Needspace{17\baselineskip}
\section{挠率、曲率与Bianchi恒等式}
\begin{definition}[挠率与曲率]\label{def:torsion-curvature}
切丛联络的挠率（torsion）为$T(X,Y)=\nabla_XY-\nabla_YX-[X,Y]$。任意向量丛联络的曲率（curvature）为
\[
 R^E(X,Y)s=\nabla_X\nabla_Ys-\nabla_Y\nabla_Xs-\nabla_{[X,Y]}s.
\]
切丛上省略上标$E$。$T=0$称无挠；$R^E=0$称联络平坦（flat）。这里平坦按曲率定义；对允许挠率的切丛，不预先把它等同于“存在坐标使全部$\Gamma$为零”。
\end{definition}

\begin{proposition}[张量性与分量]\label{prop:curvature-tensor}
$T$对两个输入均为$C^\infty(M)$线性，是$\Lambda^2T^*M\otimes TM$截面；$R^E$对$X,Y,s$均为函数线性，且对$X,Y$反对称，所以$R^E\in\Omega^2(M;\operatorname{End}E)$。坐标切标架中
\begin{align*}
 T^k{}_{ij}&=\Gamma^k_{ij}-\Gamma^k_{ji},\\
 R(\partial_i,\partial_j)\partial_k&=R^\ell{}_{kij}\partial_\ell,\\
 R^\ell{}_{kij}&=\partial_i\Gamma^\ell_{jk}-\partial_j\Gamma^\ell_{ik}
 +\Gamma^\ell_{im}\Gamma^m_{jk}-\Gamma^\ell_{jm}\Gamma^m_{ik}.
\end{align*}
\end{proposition}
\begin{proof}
例如$[fX,Y]=f[X,Y]-Y(f)X$，恰抵消$\nabla_Y(fX)$中的额外项，故$T(fX,Y)=fT(X,Y)$；反对称性给另一输入的线性。

对曲率的截面输入展开两重Leibniz，第一项给
\[
 \nabla_X\nabla_Y(fs)=f\nabla_X\nabla_Ys+X(f)\nabla_Ys
 +Y(f)\nabla_Xs+X(Y(f))s.
\]
减去交换项和$\nabla_{[X,Y]}(fs)$，一阶项抵消，零阶余项为$(XY-YX-[X,Y])(f)s=0$。对方向$fX$，第二项产生$-Y(f)\nabla_Xs$，括号项产生相反的一项，亦相消；另一方向用反对称性。局部系数光滑，所以这些函数线性映射确为所列光滑张量。

在坐标基中$[\partial_i,\partial_j]=0$，挠率式直接得到。计算$\nabla_{\partial_i}(\Gamma^m_{jk}\partial_m)$，分别对系数和标架求导，再减去$i,j$交换式，即得曲率分量。一般非坐标标架还须保留括号项，不能只把偏导换成方向导数就结束。
\end{proof}

\begin{proposition}[结构方程]\label{prop:cartan-structure}
矩阵值形式的乘法约定为$(\alpha\wedge\beta)^a{}_b=\sum_c\alpha^a{}_c\wedge\beta^c{}_b$。若$\omega$为$E$的联络矩阵，则曲率矩阵为
\[
 \Omega=\dd\omega+\omega\wedge\omega,
 \qquad R^E(X,Y)(e_b)=e_a\Omega^a{}_b(X,Y).
\]
换标架$e'=eG$时$\Omega'=G^{-1}\Omega G$。在切丛标架$e_a$及其对偶余标架$\theta^a$中，挠率形式为
\[
 \Theta^a=\dd\theta^a+\omega^a{}_b\wedge\theta^b,
 \qquad T(X,Y)=e_a\Theta^a(X,Y).
\]
若$[e_i,e_j]=c^k{}_{ij}e_k$且$\nabla_{e_i}e_j=\Gamma^a_{ij}e_a$，则$T^a{}_{ij}=\Gamma^a_{ij}-\Gamma^a_{ji}-c^a{}_{ij}$。
\end{proposition}
\begin{proof}
对局部系数列向量$s$，算子$\nabla_X=X+\omega(X)$。两次复合后，作用在$s$上的一阶项抵消，余下矩阵是
\[
 X(\omega(Y))-Y(\omega(X))-\omega([X,Y])
 +\omega(X)\omega(Y)-\omega(Y)\omega(X),
\]
恰为$\Omega(X,Y)$。曲率已经证明为内禀的端同态，因此将输入、输出标架同时换为$e'$，矩阵自然为$G^{-1}\Omega G$；这也可将$\omega'$代入验证。

对切丛，先在标架$(e_i,e_j)$求值：$\dd\theta^a(e_i,e_j)=-c^a{}_{ij}$，而$(\omega^a{}_b\wedge\theta^b)(e_i,e_j)=\Gamma^a_{ij}-\Gamma^a_{ji}$。和为$T$的分量；两端均为张量，故对任意$X,Y$成立。对偶列$\theta$换标架为$G^{-1}\theta$，相应$\Theta'=G^{-1}\Theta$，表达式与几何对象相容。
\end{proof}

\Needspace{17\baselineskip}
\begin{theorem}[一般联络的Bianchi恒等式]\label{thm:bianchi-forms}
切丛联络的挠率与曲率满足第一结构形式的Bianchi恒等式
\[
 D\Theta:=\dd\Theta+\omega\wedge\Theta=\Omega\wedge\theta.
\]
任意向量丛联络的曲率满足第二Bianchi恒等式
\[
 D\Omega:=\dd\Omega+\omega\wedge\Omega-\Omega\wedge\omega=0.
\]
矩阵乘法次序按定义保留，无须度规或无挠假设。若$T=0$，第一式推出$\sum_{\mathrm{cyc}}R(X,Y)Z=0$；有挠时不能直接删去左端。协变张量形式及挠率修正将在下一节完整写出。
\end{theorem}
\begin{proof}
由$\Theta=\dd\theta+\omega\wedge\theta$和次数为一的乘法法则，
\begin{align*}
 \dd\Theta+\omega\wedge\Theta
 &=\dd\omega\wedge\theta-\omega\wedge\dd\theta
   +\omega\wedge\dd\theta+\omega\wedge\omega\wedge\theta\\
 &=(\dd\omega+\omega\wedge\omega)\wedge\theta=\Omega\wedge\theta.
\end{align*}
对曲率有$\dd\Omega=\dd\omega\wedge\omega-\omega\wedge\dd\omega$，所以
\begin{align*}
 \dd\Omega+\omega\wedge\Omega-\Omega\wedge\omega
 &=\dd\omega\wedge\omega-\omega\wedge\dd\omega
   +\omega\wedge(\dd\omega+\omega\wedge\omega)\\
 &\quad-(\dd\omega+\omega\wedge\omega)\wedge\omega=0.
\end{align*}
这是结合律下的逐项抵消，不曾交换矩阵次序。$T=0$时$\Theta=0$，在$(X,Y,Z)$上求$\Omega\wedge\theta$的值，即为$R(X,Y)Z+R(Y,Z)X+R(Z,X)Y$，从而得到循环恒等式。
\end{proof}

\begin{theorem}[平坦性与局部平行标架]\label{thm:flat-parallel-frame}
向量丛联络满足$R^E=0$，当且仅当每点附近存在平行标架，即$\nabla_Xe_a=0$对所有$X,a$成立。在这种邻域内，平行移动只依赖端点。对切丛，进一步有$T=0$时才可由此取局部坐标，使坐标标架平行、全部$\Gamma$为零。
\end{theorem}
\begin{proof}
平行标架中$\omega=0$，结构方程立刻给$\Omega=0$。反向，在坐标盒内把$E$平凡化，写$\omega=A_i(x)\dd x^i$。在总空间上定义水平向量场
\[
 H_i=\partial_{x^i}-(A_i(x)v)^a\partial_{v^a}.
\]
直接对坐标函数求括号可得$[H_i,H_j]=-(\Omega_{ij}(x)v)^a\partial_{v^a}$，因此平坦时各$H_i$交换。它们的流在底空间平移相应坐标，纤维部分解线性方程；将底空间盒缩小后，有限段流对任意初始纤维向量都存在。由第5章“括号为零则局部流交换”的定理，从盒中心沿各坐标方向依次移动，得到与次序无关的映射$(x,v_0)\mapsto(x,P(x)v_0)$。线性ODE使$P(x)$线性且可逆，并光滑依赖$x$；交换流给$\partial_iP=-A_iP$。于是$P$的各列定义$r$个线性无关平行截面。

平行标架中沿任意曲线的系数恒定，所以该盒内的移动只取决于端点。对切丛平行标架，有$T(e_i,e_j)=-[e_i,e_j]$；无挠时标架交换，以交换流构造坐标，其坐标向量就是$e_i$。反向，坐标标架平行必同时使$T,R$为零。仅曲率为零时可能存在不交换的平行标架，不能据此断言有平行坐标基。
\end{proof}

\begin{example}[一点的系数与整个邻域的平坦性]
在$\R^2$取唯一非零系数$\Gamma^2_{22}=x^1$。它无挠，且原点所有系数为零，但$R(\partial_1,\partial_2)\partial_2=\partial_2$。因此在一点消去系数不等于在邻域内消去曲率。另一方面，平坦只保证局部平行标架；穿孔平面上的平坦线丛联络仍可有非平凡绕圈移动，见习题10.5。
\end{example}
\begin{figure}[htbp]
\centering\DGDiagram{connection-loop}

\caption{常矩阵联络$\omega=A\dd x+B\dd y$沿矩形的移动次序。按$+x,+y,-x,-y$走，作用于列向量的总矩阵从右往左依次相乘，其首个非平凡项为$-\varepsilon\delta[A,B]$。这是明确的矩阵模型，不以图替代一般曲率定理。原创示意。}
\label{fig:connection-loop}
\end{figure}


\Needspace{17\baselineskip}
\section{张量丛上的诱导联络}
\begin{theorem}[对偶、张量积及外幂联络]\label{thm:induced-connection}
给定$E,F$上的联络，存在唯一的对偶及张量积联络，使
\begin{align*}
 X\bigl(\alpha(s)\bigr)&=(\nabla_X\alpha)(s)+\alpha(\nabla_Xs),\\
 \nabla_X(s\otimes t)&=(\nabla_Xs)\otimes t+s\otimes\nabla_Xt.
\end{align*}
直和按分量定义。诱导联络保持置换、缩并与交替化，从而在所有张量丛和外幂丛上给出联络。在对偶标架中，矩阵为$-\omega^{\mathsf T}$。
\end{theorem}
\begin{proof}
由第一式只能定义$(\nabla_X\alpha)(s)=X(\alpha(s))-\alpha(\nabla_Xs)$。将$s$换成$fs$，两个$X(f)\alpha(s)$抵消，所以右端对$s$函数线性，确为光滑对偶截面。在局部对偶基$e^a$上由$e^a(e_b)=\delta^a_b$得$\nabla_Xe^a=-\omega^a{}_b(X)e^b$，即列系数作用矩阵为$-\omega^{\mathsf T}$，并直接满足联络公理。

张量积按第二式定义时，关系$(fs)\otimes t=s\otimes(ft)$两边求导都产生一次$X(f)s\otimes t$及相同的其余项，所以与张量积的平衡关系相容。纯张量张成局部截面，故定义唯一，局部展开也保证光滑和两条函数规则。逐因子求导与置换交换，与自然求值配对相容；因此与缩并及交替投影相容，可限制到对称或交替子丛。没有在不同坐标中另行选择不相容的导数。
\end{proof}

\Needspace{20\baselineskip}
\begin{proposition}[混合张量的坐标公式]\label{prop:tensor-covariant-derivative}
对仿射联络及$(r,s)$型张量场$Q$，协变导数是$(r,s+1)$型张量，方向槽放在最后；在坐标中记$\nabla_iQ=\nabla_{\partial_i}Q$，分量为
\begin{align*}
 (\nabla_iQ)^{a_1\cdots a_r}{}_{b_1\cdots b_s}
 &=\partial_iQ^{a_1\cdots a_r}{}_{b_1\cdots b_s}\\
 &\quad+\sum_{u=1}^r\Gamma^{a_u}_{ic}
 Q^{a_1\cdots c\cdots a_r}{}_{b_1\cdots b_s}\\
 &\quad-\sum_{v=1}^s\Gamma^c_{ib_v}
 Q^{a_1\cdots a_r}{}_{b_1\cdots c\cdots b_s}.
\end{align*}
特别地$\nabla_if=\partial_if$、$\nabla_i\alpha_j=\partial_i\alpha_j-\Gamma^k_{ij}\alpha_k$。对任意一形式，
\[
 \dd\alpha(X,Y)=(\nabla_X\alpha)(Y)-(\nabla_Y\alpha)(X)+\alpha(T(X,Y)).
\]
\end{proposition}
\begin{proof}
在局部张量标架上逐因子使用上一条定理：向量基贡献正的$\Gamma$，对偶基贡献负的$\Gamma$，系数贡献偏导，合并即得公式。由于定义内禀且对方向函数线性，新添槽为协变槽；上式中$i$写在前面只是便于计算，不改变型别或槽位约定。

由$\dd\alpha(X,Y)=X(\alpha(Y))-Y(\alpha(X))-\alpha([X,Y])$，把前两项分别换成协变导数加上$\alpha(\nabla_XY)$和$\alpha(\nabla_YX)$，余下恰为$\alpha(T(X,Y))$。因此只有无挠时才可直接把外微分写成协变导数的反对称化。
\end{proof}

\begin{proposition}[诱导曲率与二阶导数]\label{prop:ricci-commutator}
对偶曲率满足$(R^{E^*}(X,Y)\alpha)(s)=-\alpha(R^E(X,Y)s)$，张量积曲率为$R^{E\otimes F}=R^E\otimes I+I\otimes R^F$。所以混合张量的曲率在每个上指标上按正号作用，在每个下指标上按负号作用。

若另有底流形的仿射联络，定义沿$X,Y$的协变Hessian
\[
 (\nabla^2s)(X,Y)=\nabla_X\nabla_Ys-\nabla_{\nabla_XY}s.
\]
则
\[
 (\nabla^2s)(X,Y)-(\nabla^2s)(Y,X)
 =R^E(X,Y)s-\nabla_{T(X,Y)}s.
\]
对函数$f$即为$\operatorname{Hess}f(X,Y)-\operatorname{Hess}f(Y,X)=-\dd f(T(X,Y))$。
\end{proposition}
\begin{proof}
在配对$\alpha(s)$上施加$XY-YX-[X,Y]=0$。用两次配对乘法法则展开，含一次导数的交叉项抵消，只剩$(R^{E^*}\alpha)(s)+\alpha(R^Es)=0$，得到负号。对$s\otimes t$同样展开两次张量积乘法法则，交叉项成对抵消，只保留各因子的曲率。由逐因子构造得到一般张量结论。

Hessian差按定义为两重导数的交换子减去$\nabla_{\nabla_XY-\nabla_YX}$；用$\nabla_XY-\nabla_YX=[X,Y]+T(X,Y)$整理即得。对函数的曲率算子为零，因为$XYf-YXf-[X,Y]f=0$，但完整Hessian还对新增的协变槽求导，所以有挠时它并不对称。不能把$\partial_i\partial_jf$与完整协变二阶导数混为一谈。
\end{proof}

\begin{theorem}[含挠率的张量形式Bianchi恒等式]\label{thm:bianchi-tensors}
令$\sum_{\mathrm{cyc}}$按$(X,Y,Z)$的三次循环求和。用上面诱导的张量联络，任意仿射联络满足
\begin{align*}
 \sum_{\mathrm{cyc}}R(X,Y)Z
 &=\sum_{\mathrm{cyc}}\bigl((\nabla_XT)(Y,Z)+T(T(X,Y),Z)\bigr),\\
 \sum_{\mathrm{cyc}}\bigl((\nabla_XR)(Y,Z)+R(T(X,Y),Z)\bigr)&=0.
\end{align*}
第二式为端同态恒等式，可再作用于任意$W$。$T=0$时它们化为经典的两条循环Bianchi恒等式。
\end{theorem}
\begin{proof}
先明确求导槽：例如
\begin{align*}
 (\nabla_XR)(Y,Z)W
 &=\nabla_X(R(Y,Z)W)-R(\nabla_XY,Z)W\\
 &\quad-R(Y,\nabla_XZ)W-R(Y,Z)\nabla_XW.
\end{align*}
$\nabla_XT$也对两个输入各扣除一项。对$TM$值二形式$\Theta$，在三向量上展开$D\Theta$，得到
\[
 (D\Theta)(X,Y,Z)=\sum_{\mathrm{cyc}}
 \bigl(\nabla_X(T(Y,Z))-T([X,Y],Z)\bigr).
\]
将第一项写成
\[
 \nabla_X(T(Y,Z))=(\nabla_XT)(Y,Z)+T(\nabla_XY,Z)+T(Y,\nabla_XZ).
\]
循环和中后两类合并为$T(\nabla_XY-\nabla_YX,Z)$；再减去括号项，恰为$T(T(X,Y),Z)$。上一节$D\Theta=\Omega\wedge\theta$的右端是循环$R(X,Y)Z$，给出第一式。

对端同态值二形式$\Omega$，$D\Omega=0$求值为
\[
 \sum_{\mathrm{cyc}}\bigl(\nabla_X^{\operatorname{End}}(R(Y,Z))-R([X,Y],Z)\bigr)=0.
\]
将端同态导数分成完整$\nabla_XR$与两个输入导数；同样把这些输入项按循环合并，得到$R(T(X,Y),Z)$。这正是第二式。证明只用了已建立的结构方程和诱导联络，没有预设无挠或度规曲率对称性。
\end{proof}

\begin{example}[平行移动与张量运算相容]
若$V(t)$和$\alpha(t)$分别是平行的向量与对偶向量，则$\frac{\dd}{\dd t}\alpha(V)=(D_t\alpha)(V)+\alpha(D_tV)=0$。因此对偶移动是$P^{-1}$的对偶，张量积移动逐因子作用。并无一般结论说$g(V,V)$恒定：这还须$\nabla g=0$；下一章才研究由度规唯一确定的无挠联络。
\end{example}


\Needspace{15\baselineskip}
\section{联络与曲率习题及详细解答}
习题10.1改述梁书上册第二版第3章原题2，10.2改述并扩展原题1，原题均在印刷页84/PDF100；10.3合并原题6与8的沿曲线公式和移动可逆任务，原题在印刷页85/PDF101。其余5题为原创补充。所有符号按本书“方向为第一个下指标”的约定重写，全部解答由本工程编写。

来源分类与原题号见\texttt{references/EXERCISES.md}。

\Needspace{12\baselineskip}
\subsection*{习题10.1：偏导分量为何不是张量}
改述梁书第3章原题2。设$Y$为向量场，$A^k{}_i=\partial_iY^k$。推导坐标变化下$A$的实际变换律；再在$x>0$上取$y=x^2$、$Y=\partial_x$，给出不服从$(1,1)$型张量变换律的具体数值。
\begin{proof}[解答]
记$J^a{}_k=\partial y^a/\partial x^k$，$Y'^a=J^a{}_kY^k$。求$y$偏导得到
\[
 A'^a{}_b=\frac{\partial x^i}{\partial y^b}J^a{}_k A^k{}_i
 +\frac{\partial x^i}{\partial y^b}(\partial_iJ^a{}_k)Y^k.
\]
第二项含坐标变换二阶导数，张量变换律没有此项。在给定一维例子，$Y^x=1$，所以$A^x{}_x=0$；但$Y^y=2x$，$A'^y{}_y=(2x)^{-1}\partial_x(2x)=1/x$，例如$x=1$时为$1$，而零张量的任意坐标分量都应为零。标准联络的$y$坐标系数为$\Gamma'^y_{yy}=-1/(2x^2)$，恰使$\nabla_yY^y=1/x+\Gamma'^y_{yy}(2x)=0$，说明联络修正如何恢复张量性。
\end{proof}

\Needspace{12\baselineskip}
\subsection*{习题10.2：挠率与标量Hessian}
改述并扩展梁书第3章原题1：证明$T(X,Y)$为张量，且
\[
 \operatorname{Hess}f(X,Y)-\operatorname{Hess}f(Y,X)=-\dd f(T(X,Y)).
\]
在平面取唯一非零系数$\Gamma^1_{12}=1$、$f=x^1$，分别计算两侧，解释这为何不与混合偏导交换矛盾。
\begin{proof}[解答]
代入$T(fX,Y)$，$\nabla_Y(fX)$的$Y(f)X$与括号$[fX,Y]$中的$-Y(f)X$在挠率定义中相消，得到$fT(X,Y)$；另一槽由反对称性得到。因此$T$是函数双线性的张量。
Hessian定义给$X(Yf)-\dd f(\nabla_XY)$，交换后相减，再用$XYf-YXf=[X,Y]f$，得到所列负号。本例$T(\partial_1,\partial_2)=\partial_1$，故右侧为$-1$。坐标Hessian为$\partial_i\partial_jf-\Gamma^k_{ij}\partial_kf$，于是$\operatorname{Hess}_{12}f=-1$、$\operatorname{Hess}_{21}f=0$，左侧也为$-1$。两个普通二阶偏导仍都为零；不对称来自联络对协变槽的修正，而非违反分析中的偏导交换。
\end{proof}

\Needspace{12\baselineskip}
\subsection*{习题10.3：沿曲线的公式及逆移动}
合并梁书第3章原题6、8：在坐标曲线$x(t)$上推导$(D_tV)^k=\dot v^k+\Gamma^k_{ij}\dot x^iv^j$，并用矩阵方程证明平行移动是线性同构。说明对偶向量的移动如何由它决定。
\begin{proof}[解答]
在定义\ref{def:derivative-along-curve}中取$e_k=\partial_k$、$\omega^k{}_j=\Gamma^k_{ij}\dd x^i$，直接求值于$\dot x$即得公式。令$A^k{}_j(t)=\Gamma^k_{ij}(x(t))\dot x^i$，基本解矩阵满足$\dot P=-AP$、$P(s)=I$，任意初值解为$P(t)v$。再解$\dot Q=QA$、$Q(s)=I$，则$\frac{\dd}{\dd t}(QP)=QAP-QAP=0$，故$QP=I$；有限维方阵因此可逆，$Q=P^{-1}$。若跨过多个坐标块，逐段移动是有限个可逆线性映射的复合，仍可逆。

对偶移动须保持求值：若初始协向量为$\alpha$，终点应为$\alpha\circ P^{-1}$。它与任意平行向量的配对恒定，且由对偶联络的平行方程唯一确定。这是逆矩阵的转置在对偶系数列上的作用，不是直接再乘$P$。
\end{proof}

\Needspace{12\baselineskip}
\subsection*{习题10.4：极坐标的平坦联络与测地线}
用平面标准联络的极坐标系数，验证$R(\partial_r,\partial_\theta)\partial_r=R(\partial_r,\partial_\theta)\partial_\theta=0$，并写出测地线方程。普通圆周$r=R,\theta=t$是否为测地线？
\begin{proof}[解答]
已知$\nabla_r\partial_r=0$、$\nabla_r\partial_\theta=\nabla_\theta\partial_r=r^{-1}\partial_\theta$、$\nabla_\theta\partial_\theta=-r\partial_r$，且坐标括号为零。于是
\begin{align*}
 R(\partial_r,\partial_\theta)\partial_r
 &=\nabla_r(r^{-1}\partial_\theta)-\nabla_\theta0
 =-r^{-2}\partial_\theta+r^{-2}\partial_\theta=0,\\
 R(\partial_r,\partial_\theta)\partial_\theta
 &=\nabla_r(-r\partial_r)-\nabla_\theta(r^{-1}\partial_\theta)
 =-\partial_r+\partial_r=0.
\end{align*}
张量性及反对称性说明曲率全为零。测地线方程为$\ddot r-r\dot\theta^2=0$、$\ddot\theta+2\dot r\dot\theta/r=0$。给定圆周使第一式左端为$-R\ne0$，所以不是平面联络的测地线。对笛卡儿位置求两次导数，沿径向和单位角向的分量正好为$\ddot r-r\dot\theta^2$及$r\ddot\theta+2\dot r\dot\theta$，所以两式等价于笛卡儿直线匀速运动；这里只在不碰到极坐标奇点的卡中使用。
\end{proof}

\Needspace{12\baselineskip}
\subsection*{习题10.5：平坦但绕圈移动非平凡}
在穿孔平面的平凡实线丛上，取$\nabla=\dd+a\eta$，$a\in\R$，$\eta=(-y\dd x+x\dd y)/(x^2+y^2)$。求曲率和逆时针绕一圈的平行移动。$a\ne0$时是否有处处非零的全局平行截面？
\begin{proof}[解答]
线丛矩阵为标量一形式，$\eta\wedge\eta=0$且$\dd\eta=0$，故$\Omega=a\dd\eta=0$。沿单位圆参数$t\in[0,2\pi]$，$\eta(\dot\gamma)=1$，平行方程为$\dot v+av=0$，所以$P=e^{-2\pi a}$。$a\ne0$时它不是恒等映射。若有处处非零的全局平行截面，将其沿此闭曲线限制，起终点是同一截面值，却又须乘$e^{-2\pi a}$，矛盾。局部角卡中可取平行截面$e^{-a\theta}$，但绕圈后无法单值拼接。平坦定理只保证局部端点无关性，未抹去全局拓扑的绕圈效应。
\end{proof}

\Needspace{12\baselineskip}
\subsection*{习题10.6：非交换矩阵联络的曲率}
在平面平凡秩二丛上令$\omega=A\dd x+B\dd y$，$A=\left(\begin{smallmatrix}0&1\\0&0\end{smallmatrix}\right)$、$B=\left(\begin{smallmatrix}0&0\\1&0\end{smallmatrix}\right)$。求曲率，并按$+x,+y,-x,-y$走边长$\varepsilon,\delta$的矩形，算出移动矩阵的二阶主项。
\begin{proof}[解答]
系数常数使$\dd\omega=0$，但$\omega\wedge\omega=(AB-BA)\dd x\wedge\dd y$，其中$[A,B]=\operatorname{diag}(1,-1)$，故曲率非零。每段沿线性方程移动，作用于列向量的总矩阵为
\[
 P=e^{\delta B}e^{\varepsilon A}e^{-\delta B}e^{-\varepsilon A}.
\]
因为$A^2=B^2=0$，每个指数都等于$I$加相应一次项，直接相乘得
\[
 P=\begin{pmatrix}1-\varepsilon\delta&\varepsilon^2\delta\\
 -\varepsilon\delta^2&1+\varepsilon\delta+\varepsilon^2\delta^2\end{pmatrix}
 =I-\varepsilon\delta[A,B]+O((|\varepsilon|+|\delta|)^3).
\]
例如$\varepsilon=\delta=1$时$P=\left(\begin{smallmatrix}0&1\\-1&3\end{smallmatrix}\right)$，仍可逆而不为恒等。这里的负号由平行方程$\dot v=-\omega(\dot\gamma)v$和绕行方向共同决定，不能只看$\Omega$的符号猜测移动次序。
\end{proof}

\Needspace{12\baselineskip}
\subsection*{习题10.7：有挠时的第一Bianchi修正}
在$\R^3$的坐标标架$e_i$中，令$\nabla_1=\partial_1+A$、$\nabla_2=\partial_2+B$、$\nabla_3=\partial_3$，其中$A=E_{13}$、$B=E_{31}$，$E_{ij}$表示第$i$行第$j$列为$1$的矩阵单位。计算循环$R(e_1,e_2)e_3$，并核对含挠率的完整第一Bianchi式。
\begin{proof}[解答]
只有方向对$(e_1,e_2)$及其反向的曲率非零，$R(e_1,e_2)=[A,B]=E_{11}-E_{33}$，所以
\[
 R(e_1,e_2)e_3+R(e_2,e_3)e_1+R(e_3,e_1)e_2=-e_3.
\]
挠率为$T_{12}=-e_3,T_{23}=0,T_{31}=-e_1$。尽管它们的坐标系数恒定，协变导数并不为零；按完整输入槽求导，
\[
 (\nabla_1T)(e_2,e_3)=-e_3,\qquad
 (\nabla_2T)(e_3,e_1)=-e_3,\qquad
 (\nabla_3T)(e_1,e_2)=0.
\]
例如第一项来自$-T(e_2,Ae_3)=-T(e_2,e_1)=-e_3$；第二项的输出导数是$B(-e_1)=-e_3$，其余输入项为零。再算循环$T(T(e_1,e_2),e_3)$，只有$T(-e_1,e_2)=e_3$留下。右侧总和为$-2e_3+e_3=-e_3$，与曲率循环相同。无挠版本的右侧零在此不可使用，结构形式的恒等式却仍成立。
\end{proof}

\Needspace{12\baselineskip}
\subsection*{习题10.8：对偶曲率的负号与缩并}
对习题10.6的联络，求$R^{E^*}(\partial_x,\partial_y)e^1$，其中$e^1$为固定对偶标架首项。再证明恒等端同态$I\in\Gamma(E\otimes E^*)$平行，其诱导曲率为零。
\begin{proof}[解答]
向量曲率矩阵$C=[A,B]=\operatorname{diag}(1,-1)$，对偶曲率矩阵是$-C^{\mathsf T}$，因此首项为$-e^1$。也可逐点验证：对任意$s$，$(R^{E^*}e^1)(s)=-e^1(Cs)$，给同一结果。
端同态的诱导导数满足$(\nabla_XS)(s)=\nabla_X(Ss)-S\nabla_Xs$，令$S=I$得到零。曲率按两个因子分别作用，所以$R^{\operatorname{End}E}(X,Y)S=R^E(X,Y)S-SR^E(X,Y)$，在$I$上为零。这说明保持自然缩并并不要求底丛曲率为零；两因子的曲率在配对时带相反符号而抵消。
\end{proof}

